\FloatBarrier
\section{Esercitazioni su Dati Sbilanciati: CNN e Tomek Links}
\label{ch:dm-ex-imbalanced-cnn-tomek}

\subsection{Regole Operative e Metrica di Distanza}
\label{sec:dm-ex-imb-rules}

Negli esercizi di sottocampionamento informato (\textit{informed undersampling}) su dati sbilanciati nello spazio bidimensionale $(x_1, x_2)$ valgono due convenzioni operative fondamentali:
\begin{itemize}
    \item \textbf{Metrica di distanza:} si adotta la \textbf{distanza Euclidea}, definita per due punti $P = (x_{1P}, x_{2P})$ e $Q = (x_{1Q}, x_{2Q})$ come:
    \begin{equation}
        d(P, Q) = \sqrt{(x_{1P} - x_{1Q})^2 + (x_{2P} - x_{2Q})^2}
    \end{equation}
    \item \textbf{Regola di spareggio (\textit{tie-breaking rule}):} qualora due o pi\`u vicini si trovino esattamente alla medesima distanza minima dal punto esaminato, viene selezionato il vicino che precede gli altri in \textbf{ordine alfabetico}.
\end{itemize}

% ====================================================================
\FloatBarrier
\subsection{Esercizio 1: Condensed Nearest Neighbor (CNN)}
\label{sec:dm-ex-cnn}

\FloatBarrier
\subsubsection{Dataset Iniziale e Protocollo dell'Algoritmo}
\label{subsec:dm-ex-cnn-setup}

Si consideri un dataset bidimensionale di $7$ istanze ($A, B, C, D, E, F, G$), in cui la \textbf{Classe 0} costituisce la classe maggioritaria ($5$ punti) e la \textbf{Classe 1} la classe minoritaria ($2$ punti), come riportato nella \autoref{tab:dm-ex-cnn-data}.

\begin{table}[H]
\centering
\caption{Dataset bidimensionale di partenza per l'esercizio sul Condensed Nearest Neighbor (CNN).}
\label{tab:dm-ex-cnn-data}
\small
\begin{tabular}{cccc}
\toprule
\textbf{Point} & $x_1$ & $x_2$ & \textbf{Class} \\
\midrule
$A$ & $0$ & $0$ & $0$ \\
$B$ & $1$ & $0$ & $0$ \\
$C$ & $0$ & $1$ & $0$ \\
$D$ & $2$ & $1$ & $0$ \\
$E$ & $3$ & $2$ & $0$ \\
$F$ & $4$ & $0$ & $1$ \\
$G$ & $4$ & $2$ & $1$ \\
\bottomrule
\end{tabular}
\end{table}

Si applica l'algoritmo \textbf{Condensed Nearest Neighbor (CNN)}~\cite{hart1968} basato su un classificatore $1$-NN secondo il seguente protocollo:
\begin{enumerate}
    \item L'insieme condensato iniziale \`e $S = \{A, F, G\}$.
    \item I punti rimanenti fuori da $S$ vengono esaminati nell'ordine prestabilito: $B, C, D, E$.
    \item Il classificatore $1$-NN predice per ciascun punto la classe del vicino pi\`u prossimo attualmente presente in $S$: se il punto \`e classificato correttamente rimane fuori da $S$; se invece viene \textbf{classificato erroneamente} (\textit{misclassified}), viene \textbf{aggiunto immediatamente a $S$}.
    \item Si eseguono passate complete sui punti residui finch\'e un'intera passata non aggiunge alcun nuovo punto a $S$.
\end{enumerate}

\FloatBarrier
\subsubsection{Svolgimento Passo-Passo e Risultato Finale}
\label{subsec:dm-ex-cnn-steps}

\paragraph{Prima passata ($S = \{A, F, G\}$).}
\begin{enumerate}
    \item \textbf{Punto $B = (1, 0)$ (Classe reale: $0$):} $d(B, A) = 1$, $d(B, F) = 3$, $d(B, G) = \sqrt{13}$. Il vicino pi\`u prossimo in $S$ \`e $A$ ($d = 1$, classe $0$). Classe predetta $= 0 =$ reale $\implies$ \textbf{Non aggiungere}.
    \item \textbf{Punto $C = (0, 1)$ (Classe reale: $0$):} $d(C, A) = 1$, $d(C, F) = \sqrt{17}$, $d(C, G) = \sqrt{17}$. Il vicino pi\`u prossimo \`e $A$ ($d = 1$, classe $0$). Classe predetta $= 0 =$ reale $\implies$ \textbf{Non aggiungere}.
    \item \textbf{Punto $D = (2, 1)$ (Classe reale: $0$):}
    \begin{align*}
        d(D, A) &= \sqrt{2^2 + 1^2} = \sqrt{5}, \qquad
        d(D, F) = \sqrt{(-2)^2 + 1^2} = \sqrt{5}, \\
        d(D, G) &= \sqrt{(-2)^2 + (-1)^2} = \sqrt{5}
    \end{align*}
    Si verifica un pareggio a tre vie ($d = \sqrt{5}$ tra $A, F, G$). Per la regola di spareggio alfabetico si seleziona $A$ (classe $0$). Classe predetta $= 0 =$ reale $\implies$ \textbf{Non aggiungere}.
    \item \textbf{Punto $E = (3, 2)$ (Classe reale: $0$):} $d(E, G) = 1$, $d(E, F) = \sqrt{5}$, $d(E, A) = \sqrt{13}$. Il vicino pi\`u prossimo \`e $G$ ($d = 1$, classe $1$). Poich\'e $1 \neq 0$, $E$ \`e classificato erroneamente e viene \textbf{aggiunto immediatamente a $S$}: ora $S = \{A, F, G, E\}$.
\end{enumerate}

\paragraph{Seconda passata ($S = \{A, F, G, E\}$).}
Si riesaminano i punti residui $B, C, D$:
\begin{enumerate}
    \item \textbf{Punto $B = (1, 0)$:} vicino pi\`u prossimo $A$ ($d = 1$, classe $0$) $\implies$ \textbf{Non aggiungere}.
    \item \textbf{Punto $C = (0, 1)$:} vicino pi\`u prossimo $A$ ($d = 1$, classe $0$) $\implies$ \textbf{Non aggiungere}.
    \item \textbf{Punto $D = (2, 1)$:} ora $d(D, E) = \sqrt{(2-3)^2 + (1-2)^2} = \sqrt{2} < \sqrt{5}$. Il vicino pi\`u prossimo \`e $E$ ($d = \sqrt{2}$, classe $0$) $\implies$ Classe predetta $= 0 =$ reale $\implies$ \textbf{Non aggiungere}.
\end{enumerate}

La \autoref{tab:dm-ex-cnn-trace} sintetizza l'intera traccia di esecuzione. Poich\'e la seconda passata non aggiunge alcun punto, l'algoritmo termina restituendo l'insieme condensato finale $S_{\text{final}} = \{A, E, F, G\}$ (\autoref{fig:dm-ex-cnn-plot}): i punti interni ridondanti $B, C, D$ della classe maggioritaria vengono eliminati, preservando il punto di frontiera $E$.

\begin{table}[H]
\centering
\caption{Traccia delle iterazioni dell'algoritmo CNN ($S_{\text{iniziale}} = \{A, F, G\}$).}
\label{tab:dm-ex-cnn-trace}
\small
\begin{tabular}{cccccc}
\toprule
\textbf{Pass} & \textbf{Point} & \textbf{Nearest in current $S$} & \textbf{Predicted} & \textbf{True} & \textbf{Action} \\
\midrule
1 & $B$ & $A$ ($d=1$) & $0$ & $0$ & Do not add \\
1 & $C$ & $A$ ($d=1$) & $0$ & $0$ & Do not add \\
1 & $D$ & $A$ ($d=\sqrt{5}$; tie con $F, G$) & $0$ & $0$ & Do not add \\
1 & $E$ & $G$ ($d=1$) & $1$ & $0$ & \textbf{Add $E$} \\
2 & $B$ & $A$ ($d=1$) & $0$ & $0$ & Do not add \\
2 & $C$ & $A$ ($d=1$) & $0$ & $0$ & Do not add \\
2 & $D$ & $E$ ($d=\sqrt{2}$) & $0$ & $0$ & Do not add \\
\bottomrule
\end{tabular}
\end{table}

\begin{figure}[H]
\centering
\resizebox{\columnwidth}{!}{%
\begin{tikzpicture}[scale=0.85]
    % Prima di CNN
    \begin{scope}[shift={(0,0)}]
        \draw[step=1cm, gray!25, very thin] (-0.4,-0.4) grid (5.0,2.8);
        \draw[->, gray!70] (-0.4,0) -- (5.2,0) node[right, font=\scriptsize] {$x_1$};
        \draw[->, gray!70] (0,-0.4) -- (0,3.0) node[above, font=\scriptsize] {$x_2$};
        % Classe 0
        \foreach \x/\y/\lbl in {0/0/A, 1/0/B, 0/1/C, 2/1/D, 3/2/E} {
            \fill[navy] (\x,\y) circle (4pt);
            \node[above right, font=\scriptsize\bfseries, color=navy] at (\x,\y) {\lbl};
        }
        % Classe 1
        \foreach \x/\y/\lbl in {4/0/F, 4/2/G} {
            \node[regular polygon, regular polygon sides=3, fill=crimson, inner sep=2.2pt] at (\x,\y) {};
            \node[above right, font=\scriptsize\bfseries, color=crimson] at (\x,\y) {\lbl};
        }
        \node[font=\footnotesize\bfseries] at (2.3,-0.8) {(a) Dataset originale};
    \end{scope}

    % Dopo CNN
    \begin{scope}[shift={(7.8,0)}]
        \draw[step=1cm, gray!25, very thin] (-0.4,-0.4) grid (5.0,2.8);
        \draw[->, gray!70] (-0.4,0) -- (5.2,0) node[right, font=\scriptsize] {$x_1$};
        \draw[->, gray!70] (0,-0.4) -- (0,3.0) node[above, font=\scriptsize] {$x_2$};
        % Punti rimossi tratteggiati
        \foreach \x/\y/\lbl in {1/0/B, 0/1/C, 2/1/D} {
            \draw[gray!60, dashed] (\x,\y) circle (4pt);
            \node[above right, font=\scriptsize, color=gray] at (\x,\y) {\lbl};
        }
        % Classe 0 superstiti
        \foreach \x/\y/\lbl in {0/0/A, 3/2/E} {
            \fill[navy] (\x,\y) circle (4pt);
            \node[above right, font=\scriptsize\bfseries, color=navy] at (\x,\y) {\lbl};
        }
        % Classe 1
        \foreach \x/\y/\lbl in {4/0/F, 4/2/G} {
            \node[regular polygon, regular polygon sides=3, fill=crimson, inner sep=2.2pt] at (\x,\y) {};
            \node[above right, font=\scriptsize\bfseries, color=crimson] at (\x,\y) {\lbl};
        }
        \node[font=\footnotesize\bfseries] at (2.3,-0.8) {(b) $S_{\text{final}} = \{A, E, F, G\}$};
    \end{scope}
\end{tikzpicture}%
}
\caption{Effetto dell'algoritmo CNN: (a) dataset iniziale; (b) punti superstiti dopo la rimozione delle istanze ridondanti $B, C, D$ della Classe $0$.}
\label{fig:dm-ex-cnn-plot}
\end{figure}

% ====================================================================
\FloatBarrier
\subsection{Esercizio 2: Tomek Links}
\label{sec:dm-ex-tomek}

Si consideri il dataset bidimensionale di $7$ punti riportato nella \autoref{tab:dm-ex-tomek-data}, in cui la \textbf{Classe 0} ($4$ istanze) viene sottoposta a sottocampionamento rispetto alla \textbf{Classe 1} ($3$ istanze).

\begin{definition}[Tomek Link]
Un \textbf{Tomek Link}~\cite{tomek1976} \`e una coppia di osservazioni $(x_i, x_j)$ appartenenti a \textbf{classi opposte} che risultano essere \textbf{vicini pi\`u prossimi reciproci} (\textit{mutual nearest neighbors}).
\end{definition}

\begin{table}[H]
\centering
\caption{Dataset e identificazione dei vicini pi\`u prossimi per l'esercizio sui Tomek Links.}
\label{tab:dm-ex-tomek-data}
\small
\resizebox{\columnwidth}{!}{%
\begin{tabular}{ccccccc}
\toprule
\textbf{Point} & $x_1$ & $x_2$ & \textbf{Class} & \textbf{Nearest Neighbor} & \textbf{Distanza Euclidea} & \textbf{Classi della coppia} \\
\midrule
$A$ & $0$ & $0$ & $0$ & $B$ & $d(A,B) = 1$ & $(0,0)$ --- stessa classe \\
$B$ & $1$ & $0$ & $0$ & $A$ & $d(B,A) = 1$ & $(0,0)$ --- stessa classe \\
$C$ & $3$ & $0$ & $0$ & $E$ & $d(C,E) = 1$ & $(0,1)$ --- \textbf{classi opposte} \\
$D$ & $5$ & $0$ & $0$ & $F$ & $d(D,F) = 1$ & $(0,1)$ --- \textbf{classi opposte} \\
$E$ & $3$ & $1$ & $1$ & $C$ & $d(E,C) = 1$ & $(1,0)$ --- \textbf{classi opposte} \\
$F$ & $5$ & $1$ & $1$ & $D$ & $d(F,D) = 1$ & $(1,0)$ --- \textbf{classi opposte} \\
$G$ & $8$ & $2$ & $1$ & $F$ & $d(G,F) = \sqrt{10}$ & $(1,1)$ --- non reciproco \\
\bottomrule
\end{tabular}%
}
\end{table}

Dall'analisi delle coppie reciproche nella \autoref{tab:dm-ex-tomek-data}:
\begin{itemize}
    \item $(A, B)$ sono vicini reciproci ma appartengono alla stessa classe ($0$), quindi non formano un Tomek Link.
    \item \textbf{$(C, E)$} e \textbf{$(D, F)$} sono vicini reciproci a distanza $1$ appartenenti a classi opposte, costituendo dunque i due \textbf{Tomek Links} del dataset.
\end{itemize}

Poich\'e il sottocampionamento viene applicato esclusivamente alla Classe $0$, si rimuovono i punti $C$ e $D$, ottenendo il dataset finale $\{A, B, E, F, G\}$ ($2$ istanze di Classe $0$ e $3$ di Classe $1$, \autoref{fig:dm-ex-tomek-plot}). Si noti che i Tomek Links agiscono come tecnica di pulizia dei confini decisionali e \textbf{non impongono un bilanciamento numerico esatto tra le classi}.

\begin{figure}[H]
\centering
\resizebox{\columnwidth}{!}{%
\begin{tikzpicture}[scale=0.65]
    % Prima
    \begin{scope}[shift={(0,0)}]
        \draw[step=1cm, gray!25, very thin] (-0.5,-0.5) grid (8.8,2.8);
        \draw[->, gray!70] (-0.5,0) -- (9.0,0) node[right, font=\scriptsize] {$x_1$};
        \draw[->, gray!70] (0,-0.5) -- (0,3.0) node[above, font=\scriptsize] {$x_2$};
        \foreach \x/\y/\lbl in {0/0/A, 1/0/B, 3/0/C, 5/0/D} {
            \fill[navy] (\x,\y) circle (4.5pt);
            \node[below, font=\scriptsize\bfseries, color=navy] at (\x,\y-0.1) {\lbl};
        }
        \foreach \x/\y/\lbl in {3/1/E, 5/1/F, 8/2/G} {
            \node[regular polygon, regular polygon sides=3, fill=crimson, inner sep=2.2pt] at (\x,\y) {};
            \node[above, font=\scriptsize\bfseries, color=crimson] at (\x,\y+0.12) {\lbl};
        }
        % Evidenziazione Tomek Links
        \draw[dashed, darkgreen, line width=1pt] (3,0) -- (3,1);
        \draw[dashed, darkgreen, line width=1pt] (5,0) -- (5,1);
        \node[font=\footnotesize\bfseries] at (4.2,-1.1) {(a) Tomek Links individuati: $(C,E)$ e $(D,F)$};
    \end{scope}

    % Dopo
    \begin{scope}[shift={(10.5,0)}]
        \draw[step=1cm, gray!25, very thin] (-0.5,-0.5) grid (8.8,2.8);
        \draw[->, gray!70] (-0.5,0) -- (9.0,0) node[right, font=\scriptsize] {$x_1$};
        \draw[->, gray!70] (0,-0.5) -- (0,3.0) node[above, font=\scriptsize] {$x_2$};
        \foreach \x/\y/\lbl in {3/0/C, 5/0/D} {
            \draw[gray!60, dashed] (\x,\y) circle (4.5pt);
            \node[below, font=\scriptsize, color=gray] at (\x,\y-0.1) {\lbl};
        }
        \foreach \x/\y/\lbl in {0/0/A, 1/0/B} {
            \fill[navy] (\x,\y) circle (4.5pt);
            \node[below, font=\scriptsize\bfseries, color=navy] at (\x,\y-0.1) {\lbl};
        }
        \foreach \x/\y/\lbl in {3/1/E, 5/1/F, 8/2/G} {
            \node[regular polygon, regular polygon sides=3, fill=crimson, inner sep=2.2pt] at (\x,\y) {};
            \node[above, font=\scriptsize\bfseries, color=crimson] at (\x,\y+0.12) {\lbl};
        }
        \node[font=\footnotesize\bfseries] at (4.2,-1.1) {(b) Dataset ripulito: $\{A, B, E, F, G\}$};
    \end{scope}
\end{tikzpicture}%
}
\caption{Pulizia dei confini tramite Tomek Links: (a) coppie reciproche $(C,E)$ e $(D,F)$ evidenziate in verde; (b) dataset finale dopo la rimozione di $C$ e $D$ dalla Classe $0$.}
\label{fig:dm-ex-tomek-plot}
\end{figure}
